\section{本讲的问题：为什么 Transformer 不总是 full dense compute？}
Lecture 4 把两个看似不同的方向放在一起：attention alternatives 和 mixture of experts。前者问：长上下文下，是否每个 token 都必须 attend 所有历史？后者问：每个 token 是否必须激活模型的所有参数？两者的共同目标是 selective computation：把昂贵计算只花在必要位置。本讲因此会同时追踪质量、每 token FLOPs、状态大小、路由负载和跨设备通信，而不是只比较参数量。
\begin{importantbox}{本讲主线}
Attention alternatives 选择“看哪些历史状态”；MoE 选择“用哪些专家参数”。二者都是在模型质量、训练效率、推理成本和系统复杂度之间做条件计算 tradeoff。
\end{importantbox}
